%% ---------------------------------------------------------------------------
%% conclusiones.tex
%%
%% Capítulo 5: Conclusiones y recomendaciones.
%% Estructura aprobada el 1/11/2026 y redacción en el estilo del TFG de ejemplo
%% (2/11/2026): un párrafo sobre el proyecto y el objetivo general, uno por objetivo
%% específico (en qué consistía, qué se hizo con su subsección y si se cumplió) y uno
%% de cierre con las limitaciones que anuncia las recomendaciones (5.2, corta).
%% Grado de cumplimiento indicado por el estudiante: OE1 completado y validado por
%% el asesor (su edición del 2/11/2026); OE2, OE3 y OE4 cumplidos; la utilización de
%% recursos queda como trabajo futuro.
%%
%% Convención de escritura: una oración por línea del archivo fuente.
%% ---------------------------------------------------------------------------

\chapter{Conclusiones y recomendaciones}
\label{ch:conclusiones}

\section{Conclusiones}

% Redacción en el estilo del TFG de ejemplo (2/11/2026): por objetivo, en qué
% consistía, qué se hizo (con la subsección) y si se cumplió.
El proyecto tuvo como propósito desarrollar una plataforma SLAM híbrida capaz de operar en tiempo real, combinando la visión computacional y los datos GPS de un teléfono con un mecanismo de aprendizaje por refuerzo, para incrementar la precisión y la estabilidad de la localización en recorridos viales.
Este propósito se alcanzó.
La plataforma estima la posición de un vehículo en tiempo real con la cámara, el GPS y, con la aprobación del profesor asesor, el giroscopio del teléfono, y procesó recorridos completos en vivo durante la prueba de campo.
Frente a la cámara sola, reduce el error de cientos o miles de metros a unos pocos metros (subsección~\ref{sec:res-precision}).
La mayor parte de esa mejora se debe a la fusión del GPS y del giroscopio, mientras que el aprendizaje por refuerzo aporta una mejora moderada, concentrada en los peores casos de las pérdidas del GPS.
Los cuatro objetivos específicos se cumplieron.

% OE1: validado por el asesor (2/11/2026), con los mismos valores de RT-1 y RT-2.
\textbf{El objetivo 1} consistía en sintetizar un listado de requerimientos funcionales, computacionales y temporales para la operación en tiempo real, validado mediante criterio de experto.
Se elaboró ese listado, con una forma de verificación para cada requerimiento, a partir de la literatura y de mediciones sobre el teléfono, ya que varias de sus características reales difieren de las nominales (subsección~\ref{sec:res-requerimientos}).
También se definieron las métricas de evaluación (secciones~\ref{sec:tiempo-real} y~\ref{sec:metricas}), que después se usaron en la validación experimental.
Es decir, el objetivo se completó y se validó por parte del asesor.

\textbf{El objetivo 2} consistía en diseñar una plataforma SLAM híbrida que combinara el procesamiento visual y los datos GPS para la localización en tiempo real.
Se construyeron dos programas unidos por un protocolo de red: la aplicación del teléfono, que transmite el video, el GPS y las mediciones inerciales, y el canal de procesamiento de la computadora, con una ventana que muestra el video, el estado y el mapa (subsecciones~\ref{sec:res-arquitectura}, \ref{sec:res-app} y~\ref{sec:res-interfaz}).
Sobre esa base se resolvieron la escala de la cámara, dos errores heredados de la odometría visual y la fusión con el giroscopio y el GPS (subsecciones~\ref{sec:res-escala} a~\ref{sec:res-fusion}).
La plataforma procesa los 30 cuadros por segundo en la computadora de escritorio, con un percentil 95 de latencia de 107~ms, dentro de los requerimientos temporales, y en la prueba de campo procesó en vivo cuatro sesiones en la portátil, con una latencia mayor que no afectó la precisión (subsecciones~\ref{sec:res-tiempo-real} a~\ref{sec:res-trayectorias}).
Por lo tanto, el objetivo se cumplió.

\textbf{El objetivo 3} consistía en desarrollar un mecanismo de aprendizaje por refuerzo que optimizara de forma adaptativa la fusión sensorial.
Antes de construirlo se midió dónde había margen de mejora, y se encontró que estaba en la confianza en la rapidez de la cámara durante las pérdidas del GPS (subsección~\ref{sec:res-rl}).
Se desarrolló una política con Q-learning tabular que decide esa confianza, se evaluó dejando fuera cada sesión con criterios fijados de antemano y se integró en el canal de procesamiento y en la ventana, con un costo de décimas de milisegundo por cuadro.
Las dos iteraciones cumplieron su criterio: la segunda redujo los peores errores frente al filtro convencional y, por poco, frente a un peso fijo bien elegido, aunque su aporte es moderado y no se traslada a otro dominio (subsecciones~\ref{sec:res-aporte-rl} y~\ref{sec:res-corte-rl}).
Es decir, el objetivo se cumplió, y la política funciona integrada en el canal de procesamiento y en la ventana.

\textbf{El objetivo 4} consistía en valorar experimentalmente el desempeño del sistema por medio de métricas de precisión, estabilidad y rendimiento computacional.
Se compararon cuatro modos de la plataforma sobre siete grabaciones en vehículo, con el GPS del teléfono como referencia (subsección~\ref{sec:res-precision}).
Frente a la cámara sola, que es la línea base de SLAM exclusivamente visual, la configuración completa redujo el ATE de 2087 a 2,2~m y prácticamente eliminó los saltos de la estimación; además, la fusión rechazó los saltos y las ráfagas de un GPS degradado.
En rendimiento, se midieron la latencia, la frecuencia efectiva y el tiempo de cada etapa (subsección~\ref{sec:res-tiempo-real}).
Por lo tanto, el objetivo se cumplió, aunque la utilización de recursos computacionales no se midió y queda como trabajo futuro.

Las limitaciones de la subsección~\ref{sec:res-limitaciones} no impiden extraer estas conclusiones, pero sí restringen su alcance.
La referencia de las grabaciones propias es el GPS del teléfono, por lo que las cifras miden la consistencia con un receptor de varios metros de error y no la exactitud absoluta de la plataforma; además, la escala de la cámara limita la estimación durante las pérdidas largas del GPS.
Las líneas de trabajo que se derivan de estas limitaciones se presentan en la sección~\ref{sec:recomendaciones}.

\section{Recomendaciones}
\label{sec:recomendaciones}

A partir de los resultados obtenidos, se plantean las siguientes recomendaciones para dar continuidad al trabajo realizado.

Primero, se recomienda medir la utilización de recursos computacionales de la plataforma, es decir, el uso del procesador, de la tarjeta gráfica y de la memoria, en la computadora de escritorio y en la portátil, para completar la valoración del rendimiento computacional.

Segundo, se recomienda contar con una trayectoria de referencia de mayor precisión para las grabaciones propias, por ejemplo con un receptor de posicionamiento de alta precisión instalado en el vehículo, ya que el GPS del teléfono no permite medir la exactitud absoluta de la plataforma.

Tercero, se recomienda estudiar la estimación de la escala de la cámara a rapideces altas y en escenas lejanas, ya que es la fuente de error dominante en los cortes de más de 20~s y la que define el margen que una política podría aprovechar.

% CTX §4.37 (sugerencia del asesor, 23/10/2026): dos sesiones (14:05 y 14:18), nada reentrenado.
Cuarto, se recomienda retomar como trabajo futuro el uso del acelerómetro para validar los cambios de rapidez que indica la cámara durante una pérdida del GPS.
Una medición preliminar, con solo dos sesiones y sin reentrenar la política, mostró que desconfiar de la cámara cuando el acelerómetro no ve el cambio no reduce el error, porque descartarla equivale a suponer una rapidez constante.
Usarlo al revés, para no descartar la cámara cuando el vehículo realmente acelera o frena, mejora poco, con un techo de 5 a 6~\% del percentil 90.

Quinto, se recomienda corregir el comportamiento del modo exclusivamente visual en los altos, para eliminar los saltos de la línea base, y ampliar los datos de entrenamiento de la política con más sesiones, rutas y condiciones, ya que la política entrenada no se traslada a un dominio distinto.

%%% Local Variables:
%%% mode: latex
%%% TeX-master: "main"
%%% End:
